\documentclass[12pt,a4paper]{article}

% --- PAQUETES DE FORMATO ---
\usepackage[spanish]{babel}
\usepackage{amsmath,amssymb,amsfonts}
\usepackage{geometry}
\usepackage{booktabs}
\usepackage{hyperref}

% --- CONFIGURACIÓN DE MÁRGENES ---
\geometry{
    a4paper,
    top=2.5cm,
    bottom=2.5cm,
    left=3cm,
    right=2.5cm
}

% --- CONFIGURACIÓN DE HYPERREF ---
\hypersetup{
    colorlinks=true,
    linkcolor=black,
    filecolor=black,      
    urlcolor=black,
    pdftitle={Informe 2},
}

\begin{document}

% ==========================================
% PORTADA (Fiel al diseño de la imagen)
% ==========================================
\begin{titlepage}
    \centering
    \vspace*{6cm}
    
    {\Huge \bfseries Informe 2} \\[3cm]
    
    {\Large \bfseries Autores:} \\[0.5cm]
    \begin{tabular}{ll}
        Eric Guacache & C.I: 31.788.759 \\
        Renny Andrade & C.I: 29.727.828
    \end{tabular}
    
    \vfill
\end{titlepage}

\newpage

% ==========================================
% ÍNDICE
% ==========================================
\tableofcontents
\newpage

% ==========================================
% DESARROLLO DE LAS PREGUNTAS
% ==========================================

\section{Diferencias entre Registros Temporales (\$t0--\$t9) y Registros Guardados (\$s0--\$s7)}

\subsection{Diferencias Conceptuales}
En la arquitectura MIPS32, los registros se dividen en dos categorías principales con convenciones de uso específicas[cite: 3]:

\begin{itemize}
    \item \textbf{Registros Temporales (\$t0--\$t9):}
    \begin{itemize}
        \item \textbf{Propósito:} Almacenamiento temporal de valores intermedios durante cálculos puntuales[cite: 3].
        \item \textbf{Responsabilidad:} El llamante (\textit{caller}) debe asumir que pueden modificarse tras una llamada a función[cite: 3].
        \item \textbf{Preservación:} No se garantiza que conserven su valor después de un \texttt{jal}[cite: 3].
        \item \textbf{Uso típico:} Variables auxiliares en cálculos aritméticos, índices locales dentro de bucles de una sola subrutina[cite: 3].
    \end{itemize}
    \item \textbf{Registros Guardados (\$s0--\$s7):}
    \begin{itemize}
        \item \textbf{Propósito:} Valores que deben persistir a través de llamadas a subrutinas[cite: 3].
        \item \textbf{Responsabilidad:} El llamado (\textit{callee}) debe preservarlos en la pila (\textit{stack}) antes de modificarlos y restaurarlos antes de retornar[cite: 3].
        \item \textbf{Preservación:} Deben mantener su valor antes y después de cualquier \texttt{jal}[cite: 3].
        \item \textbf{Uso típico:} Punteros a arreglos, índices de límites (\textit{left}, \textit{right}, \textit{mid}) o variables críticas que se utilizan a lo largo de llamadas recursivas o anidadas[cite: 3].
    \end{itemize}
\end{itemize}

\subsection{Aplicación en los Algoritmos (Quicksort y Mergesort)}
Dado que tanto Quicksort como Mergesort se basan en la técnica de divide y vencerás y emplean recursividad, el uso de registros guardados (\texttt{\$s0}--\texttt{\$s7}) y del stack se vuelve fundamental en ambas implementaciones[cite: 3]:
\begin{itemize}
    \item \textbf{En Mergesort:} Se utilizan registros guardados (como \texttt{\$s0}, \texttt{\$s1} y \texttt{\$s2}) para preservar de manera segura el puntero base del arreglo, el índice inicial (\textit{left}) y el índice final (\textit{right}), así como el punto medio (\textit{mid}) calculado[cite: 3]. Al realizar las dos llamadas recursivas sucesivas (\texttt{jal mergesort} para la primera mitad y \texttt{jal mergesort} para la segunda mitad), estos valores en los registros \texttt{\$s} deben permanecer intactos para calcular correctamente los rangos del paso posterior de combinación (\texttt{jal merge})[cite: 3].
    \item \textbf{En Quicksort:} De forma análoga, registros como \texttt{\$s0} y \texttt{\$s1} almacenan los límites de la partición actual[cite: 3]. Se requiere asegurar su valor en la pila antes de ejecutar la partición (\texttt{jal partition}) y la llamada recursiva al subarreglo izquierdo[cite: 3].
\end{itemize}

\subsection{Justificación del Diseño en Algoritmos Recursivos}
\begin{itemize}
    \item \textbf{Persistencia tras invocaciones:} A diferencia de algoritmos iterativos simples, el flujo de ejecución recursivo con \texttt{jal} sobrescribe rutinariamente los registros temporales[cite: 3]. Por ello, delegar variables de estado clave a \texttt{\$s0}--\texttt{\$s7} previene la corrupción de datos al regresar de una llamada recursiva[cite: 3].
    \item \textbf{Garantía del convenio ABI:} Guardar los registros \texttt{\$s} y el registro de retorno \texttt{\$ra} en la pila al inicio de la subrutina (prólogo) y restaurarlos al final (epílogo) garantiza la corrección del programa en entornos MIPS32[cite: 3].
    \item \textbf{Uso puntual de temporales:} Los registros \texttt{\$t0}--\texttt{\$t9} se reservan únicamente dentro de las rutinas auxiliares de bajo nivel (como \texttt{merge} o \texttt{partition}) para comparaciones de elementos e incrementos de índices temporales donde no se realizan nuevas llamadas anidadas[cite: 3].
\end{itemize}


\section{Diferencias entre los Registros \$a0--\$a3, \$v0--\$v1 y \$ra}

\begin{itemize}
    \item \textbf{Registros de Argumentos (\$a0--\$a3):}
    \begin{itemize}
        \item \textbf{Propósito:} Transferir parámetros a las subrutinas[cite: 3].
        \item \textbf{Aplicación práctica:} En Quicksort, se usan \texttt{\$a0} (dirección del arreglo), \texttt{\$a1} (\textit{low}) y \texttt{\$a2} (\textit{high})[cite: 3]. En Mergesort, se utiliza adicionalmente \texttt{\$a3} para pasar el punto medio (\textit{mid}) al invocar a la rutina de combinación \texttt{merge} (\texttt{jal merge})[cite: 3].
    \end{itemize}
    \item \textbf{Registros de Retorno (\$v0--\$v1):}
    \begin{itemize}
        \item \textbf{Propósito:} Devolver valores de las funciones y seleccionar códigos de servicios del sistema (\texttt{syscall})[cite: 3].
        \item \textbf{Aplicación práctica:} La función \texttt{partition} en Quicksort calcula la posición definitiva del pivote y la retorna a través de \texttt{\$v0} (\texttt{move \$v0, \$t2})[cite: 3]. En los bloques de impresión y llamadas del sistema, se carga en \texttt{\$v0} el servicio solicitado (\texttt{li \$v0, 4} para strings, \texttt{li \$v0, 1} para enteros, \texttt{li \$v0, 10} para finalizar)[cite: 3].
    \end{itemize}
    \item \textbf{Registro de Dirección de Retorno (\$ra):}
    \begin{itemize}
        \item \textbf{Propósito:} Almacenar la dirección a la que debe volver la ejecución tras un salto \texttt{jal}[cite: 3].
        \item \textbf{Aplicación práctica:} Debido a que tanto Mergesort como Quicksort realizan llamadas recursivas (múltiples \texttt{jal}), el valor de \texttt{\$ra} se sobrescribe en cada llamada anidada[cite: 3]. Por esta razón, ambas funciones reservan espacio en el stack y salvan \texttt{\$ra} explícitamente (\texttt{sw \$ra, 16(\$sp)}) antes de llamar a subrutinas y lo restauran (\texttt{lw \$ra, 16(\$sp)}) antes de ejecutar \texttt{jr \$ra}[cite: 3].
    \end{itemize}
\end{itemize}


\section{Impacto de Registros frente a Memoria en el Rendimiento}

En MIPS32, las operaciones entre registros se ejecutan en un único ciclo de reloj dentro de la ALU, mientras que las instrucciones de acceso a memoria (\texttt{lw}/\texttt{sw}) requieren acceder al subsistema de datos, aumentando la latencia y los posibles estancamientos del pipeline[cite: 3].

\begin{table}[h!]
\centering
\caption{Tipos de acceso en cada algoritmo}
\vspace{0.2cm}
\begin{tabular}{lll}
\toprule
\textbf{Operación} & \textbf{Mergesort} & \textbf{Quicksort} \\
\midrule
\textbf{Accesos a Registros} & Alto (cálculo de rangos/índices) & Alto (comparaciones e incrementos) \\
\textbf{Accesos a Memoria (E/S)} & Muy Alto (copia constante a \texttt{temp}) & Medio-Bajo (\textit{swaps} in-place) \\
\textbf{Uso de Stack (Pila)} & Alto (20 bytes por marco) & Alto (20 bytes por marco) \\
\bottomrule
\end{tabular}
\end{table}

\begin{itemize}
    \item \textbf{Caso Mergesort:} Presenta un alto costo de memoria secundaria[cite: 3]. Durante la etapa \texttt{merge}, lee los elementos de \texttt{array}, los escribe en un arreglo auxiliar \texttt{.space 100} (\texttt{temp}), y luego ejecuta el bucle \texttt{loop\_copia} para volcar de nuevo los datos a la memoria original[cite: 3]. A pesar de apoyarse en registros temporales (\texttt{\$t0}--\texttt{\$t2}) para controlar los índices $i, j, k$, la gran cantidad de lecturas (\texttt{lw \$t4, 0(\$t3)}) y escrituras (\texttt{sw \$t5, 0(\$t3)}) genera un cuello de botella por accesos a memoria de datos[cite: 3].
    \item \textbf{Caso Quicksort:} Optimiza el acceso a memoria al operar \textit{in-place}[cite: 3]. En la rutina \texttt{partition}, los elementos se comparan directamente en registros (\texttt{lw \$t5, 0(\$t4)}) y solo se modifican en memoria cuando ocurre un intercambio (\textit{swap})[cite: 3]. No requiere un búfer temporal auxiliar[cite: 3].
\end{itemize}


\section{Impacto de Estructuras de Control y Saltos}

Las bifurcaciones condicionales (\texttt{bge}, \texttt{ble}, \texttt{bgt}) provocan penalizaciones por salto (\textit{branch penalty}) en el pipeline cuando la predicción de salto falla[cite: 3].

\begin{itemize}
    \item \textbf{En Mergesort:} Los bucles como \texttt{loop\_comparar} evalúan bordes fijos (\texttt{bgt \$t0, \$a3} y \texttt{bgt \$t1, \$a2})[cite: 3]. Dado que los límites del subarreglo son predecibles, el predictor de saltos del procesador funciona de forma eficiente[cite: 3]. Sin embargo, la estructura recursiva doble (\texttt{jal mergesort}) impone una alta carga fija por la manipulación continua de punteros de la pila (\texttt{addi \$sp, \$sp, -20})[cite: 3].
    \item \textbf{En Quicksort:} La condición del bucle \texttt{part\_loop} depende del valor de los datos comparados contra el pivote (\texttt{bge \$t5, \$t1, part\_next})[cite: 3]. Dependiendo de cuán desordenado esté el arreglo original, este salto puede volverse altamente impredecible, forzando la limpieza (\textit{flush}) del pipeline repetidamente[cite: 3].
    \item \textbf{Optimizaciones:} Reorganizar instrucciones dentro de las ranuras de retardo (\textit{delay slots}) tras las instrucciones \texttt{j} o \texttt{jal} e implementar desenrollado de bucles en las etapas de copia o intercambio de elementos[cite: 3].
\end{itemize}


\section{Diferencias de Complejidad Computacional e Implicaciones}

\begin{table}[h!]
\centering
\caption{Complejidad computacional comparada}
\vspace{0.2cm}
\begin{tabular}{lllll}
\toprule
\textbf{Algoritmo} & \textbf{Mejor caso} & \textbf{Caso promedio} & \textbf{Peor caso} & \textbf{Espacio auxiliar} \\
\midrule
\textbf{Quicksort} & $\mathcal{O}(n \log n)$ & $\mathcal{O}(n \log n)$ & $\mathcal{O}(n^2)$ & $\mathcal{O}(\log n)$ (pila) \\
\textbf{Mergesort} & $\mathcal{O}(n \log n)$ & $\mathcal{O}(n \log n)$ & $\mathcal{O}(n \log n)$ & $\mathcal{O}(n)$ (pila + \texttt{temp}) \\
\bottomrule
\end{tabular}
\end{table}

\subsection{Implicaciones para la Implementación en MIPS32}
\begin{itemize}
    \item \textbf{Mergesort:}
    \begin{itemize}
        \item \textit{Garantía de Rendimiento:} Ofrece un tiempo de ejecución acotado a $\mathcal{O}(n \log n)$ incluso en el peor de los casos, lo que evita degeneraciones de tiempo en la CPU[cite: 3].
        \item \textit{Impacto en Memoria:} Requiere reservar un arreglo auxiliar explícito en la sección \texttt{.data} (\texttt{temp: .space 100})[cite: 3]. En ensamblador MIPS32, transferir continuamente los datos de \texttt{array} a \texttt{temp} y copiar el resultado ordenado de vuelta implica una cantidad masiva de instrucciones \texttt{lw} y \texttt{sw} dentro del bucle \texttt{copiar\_a\_original}, sobrecargando el bus de memoria de datos[cite: 3].
        \item \textit{Uso de Pila:} Genera marcos de pila uniformes de 20 bytes (\texttt{addi \$sp, \$sp, -20}) para almacenar \texttt{\$ra}, \texttt{\$s0} (el punto medio \textit{mid}) y los argumentos[cite: 3].
    \end{itemize}
    \item \textbf{Quicksort:}
    \begin{itemize}
        \item \textit{Sensibilidad a los Datos:} Su eficiencia depende de la elección del pivote[cite: 3]. Si el arreglo ya está ordenado y se toma el último elemento como pivote, el árbol de recursión se degrada a profundidad $\mathcal{O}(n)$, cayendo a una complejidad cuadrática $\mathcal{O}(n^2)$[cite: 3].
        \item \textit{Eficiencia In-Place:} No utiliza memoria auxiliar externa para los elementos[cite: 3]. Todas las reordenaciones se realizan intercambiando registros mediante punteros directos en la rutina \texttt{partition}[cite: 3]. Esto minimiza drásticamente los accesos al subsistema de memoria principal en comparación con Mergesort[cite: 3].
        \item \textit{Uso de Pila:} Al igual que Mergesort, reserva 20 bytes por llamada recursiva para preservar \texttt{\$ra}, \texttt{\$s0} (índice \textit{pi}) y los rangos \textit{low/high}[cite: 3].
    \end{itemize}
\end{itemize}

\subsection{Comparación Práctica y Recomendaciones}
\begin{table}[h!]
\centering
\caption{Impacto práctico en la implementación MIPS32}
\vspace{0.2cm}
\begin{tabular}{lll}
\toprule
\textbf{Factor} & \textbf{Quicksort} & \textbf{Mergesort} \\
\midrule
\textbf{Instrucciones} & Menor cantidad (uso \textit{in-place}) & Mayor cantidad (bucle de copia \texttt{merge}) \\
\textbf{Tráfico Memoria} & Moderado (solo escrituras en \textit{swaps}) & Muy alto (lecturas/escrituras en \texttt{temp}) \\
\textbf{Predictibilidad} & Variable (depende del pivote) & Alta (condiciones fijas $i, j, k$) \\
\textbf{Memoria RAM} & Arreglo original + pila & Arreglo original + Buffer \texttt{temp} \\
\bottomrule
\end{tabular}
\end{table}

\begin{itemize}
    \item \textbf{Priorizar Quicksort para sistemas restringidos:} Dado que MIPS32 es común en sistemas embebidos, la menor huella de memoria de Quicksort (\textit{in-place}) reduce la presión sobre el caché de datos y ahorra memoria principal[cite: 3].
    \item \textbf{Optimización del pivote:} Implementar la técnica de la mediana de tres para elegir el pivote garantizaría mantener la complejidad $\mathcal{O}(n \log n)$ sin caer en el peor caso[cite: 3].
    \item \textbf{Optimización de Mergesort:} Para reducir el \textit{overhead}, la etapa \texttt{copiar\_a\_original} puede optimizarse intercalando punteros o utilizando asignaciones dinámicas que eviten el volcado constante del buffer[cite: 3].
\end{itemize}


\section{Fases del Ciclo de Ejecución en MIPS32 (Camino de Datos)}

\begin{enumerate}
    \item \textbf{Búsqueda de la Instrucción (IF - Instruction Fetch):} Se utiliza el Contador de Programa (PC) para leer la dirección de la siguiente instrucción desde la memoria de instrucciones[cite: 3]. La instrucción leída se almacena en el registro de instrucción (IR) y el puntero se actualiza: $PC \leftarrow PC + 4$[cite: 3].
    \item \textbf{Decodificación y Lectura de Registros (ID - Instruction Decode):} La unidad de control interpreta el código de operación (\textit{opcode})[cite: 3]. Se leen los valores almacenados en los registros fuente ($rs$, $rt$) dentro del banco de registros[cite: 3]. Si la instrucción contiene un valor inmediato, este se extiende en signo a 32 bits[cite: 3].
    \item \textbf{Ejecución / Cálculo de Dirección (EX - Execute):} La ALU realiza el cómputo principal[cite: 3]. En operaciones Tipo R, ejecuta la operación matemática o lógica[cite: 3]. En Tipo I de memoria, suma la dirección base con el desplazamiento (\textit{offset})[cite: 3]. En saltos condicionales, evalúa la comparación[cite: 3].
    \item \textbf{Acceso a Memoria (MEM - Memory Access):} Interacción con la memoria principal de datos[cite: 3]. Para lecturas (\texttt{lw}), se extrae el dato en la dirección calculada[cite: 3]. Para escrituras (\texttt{sw}), el valor especificado se guarda en memoria[cite: 3]. Para Tipo R, esta etapa permanece inactiva[cite: 3].
    \item \textbf{Escritura de Resultados (WB - Write Back):} El valor procesado (de la ALU o de la memoria durante \texttt{lw}) se escribe permanentemente en el registro destino correspondiente ($rd$ o $rt$)[cite: 3].
\end{enumerate}


\section{Predominio de Tipos de Instrucciones (R, I, J)}

\begin{table}[h!]
\centering
\caption{Distribución estimada de tipos de instrucciones}
\vspace{0.2cm}
\begin{tabular}{llll}
\toprule
\textbf{Algoritmo} & \textbf{Tipo R} & \textbf{Tipo I} & \textbf{Tipo J} \\
\midrule
\textbf{Mergesort} & 35\% & 55\% & 10\% \\
\textbf{Quicksort} & 45\% & 50\% & 5\% \\
\bottomrule
\end{tabular}
\end{table}

\begin{itemize}
    \item \textbf{Instrucciones Tipo I (Predominantes):} Ambas implementaciones manejan el stack en MIPS32 mediante desplazamientos inmediatos (\texttt{addi \$sp, \$sp, -20}, \texttt{sw \$ra, 16(\$sp)}), calculan direcciones sumando inmediatos o constantes, y ejecutan la mayoría de saltos condicionales de control de flujo (\texttt{bge}, \texttt{bgt}, \texttt{ble})[cite: 3]. En Mergesort se observa una fuerte presencia de lecturas y escrituras directas (\texttt{lw \$t4, 0(\$t3)}, \texttt{sw \$t5, 0(\$t3)})[cite: 3].
    \item \textbf{Instrucciones Tipo R:} Se emplean para operaciones aritmético-lógicas entre registros (\texttt{add}, \texttt{sll}, \texttt{move}) y para el retorno explícito de funciones (\texttt{jr \$ra})[cite: 3].
    \item \textbf{Instrucciones Tipo J:} Representan la proporción menor[cite: 3]. Se reducen a saltos incondicionales absolutos dentro de los bucles de iteración (\texttt{j}) y a la invocación de procedimientos (\texttt{jal})[cite: 3].
\end{itemize}


\section{Impacto del Abuso de Instrucciones de Salto}

\begin{itemize}
    \item \textbf{Penalizaciones por Salto (Branch Penalty):} La arquitectura MIPS32 utiliza un pipeline de 5 etapas[cite: 3]. Si el procesador predice incorrectamente si el salto se tomará o no, las instrucciones leídas de forma especulativa deben descartarse (\textit{flush} del pipeline), lo que desperdicia ciclos de reloj[cite: 3].
    \item \textbf{Desperdicio en la Ranura de Retardo (Branch Delay Slot):} MIPS32 ejecuta por diseño la instrucción que se encuentra inmediatamente después de un salto[cite: 3]. Si no es posible reordenar el código con una instrucción útil, es necesario insertar instrucciones nulas (\texttt{nop}), aumentando el tamaño del código y consumiendo tiempo de CPU[cite: 3].
    \item \textbf{Pérdida de Localidad de Código y Misses de Caché:} Los saltos constantes a etiquetas distantes rompen la localidad espacial del código, aumentando la tasa de fallos (\textit{cache misses}) en la I-Cache y obligando a acceder a la memoria principal[cite: 3].
\end{itemize}


\section{Ventajas del Modelo RISC de MIPS}

\begin{itemize}
    \item \textbf{Conjunto de Instrucciones Reducido y Formato Fijo:} Todas las instrucciones tienen un tamaño uniforme de 32 bits, lo que simplifica la unidad de decodificación y permite que las operaciones fundamentales se ejecuten a alta velocidad en muy pocos ciclos[cite: 3].
    \item \textbf{Arquitectura Load/Store y Banco de Registros Amplio:} MIPS32 obliga a procesar datos únicamente dentro de su banco de 32 registros generales[cite: 3]. En Quicksort, esto permite comparar e intercambiar elementos en registros temporales sin saturar la memoria[cite: 3].
    \item \textbf{Segmentación de Tuberías (Pipelining) Altamente Eficiente:} El diseño uniforme de 5 etapas maximiza el rendimiento (\textit{throughput}) en bucles repetitivos de ordenamiento[cite: 3].
\end{itemize}


\section{Uso del Modo Paso a Paso (Step / Step Into) en MARS}

\begin{itemize}
    \item \textbf{Uso del botón Step (F7):} Permitió avanzar la ejecución instrucción por instrucción dentro del mismo nivel de contexto[cite: 3]. Se empleó en rutinas iterativas (\texttt{loop\_comparar}, \texttt{part\_loop}) para verificar en tiempo real cómo progresaban los índices ($i, j, k$) y si las condiciones de salida actuaban en el momento preciso[cite: 3].
    \item \textbf{Uso del botón Step Into:} Permitió "entrar" en el cuerpo de las funciones recursivas tras un \texttt{jal} para supervisar el prólogo del stack, confirmando la correcta reserva de espacio (\texttt{addi \$sp, \$sp, -20}) y el respaldo de \texttt{\$ra} y \texttt{\$s0}[cite: 3].
    \item \textbf{Validación de la Recursión:} Permitió observar la acumulación progresiva de marcos en la pila durante el proceso de división y verificar el proceso de retorno (epílogo) sin desbordar la pila[cite: 3].
\end{itemize}


\section{Herramienta de MARS para Inspección de Registros}

La herramienta más útil fue la tabla de registros (\textbf{Registers}), ubicada en el panel lateral derecho del simulador[cite: 3]:

\begin{itemize}
    \item \textbf{Monitoreo en Tiempo Real:} Muestra el valor en hexadecimal y decimal de los 32 registros, permitiendo supervisar la evolución de los índices temporales (\texttt{\$t0}--\texttt{\$t6}) y los punteros de memoria[cite: 3].
    \item \textbf{Detección Eficiente de Errores:} Facilitó verificar que los registros guardados (\texttt{\$s0}) y \texttt{\$ra} conservaran sus datos tras cada llamada recursiva, y confirmó la entrega del valor de retorno en \texttt{\$v0}[cite: 3].
    \item \textbf{Eliminación de Overhead:} Evitó la necesidad de insertar múltiples llamadas al sistema (\texttt{syscall}) para imprimir datos intermedios, manteniendo el código limpio[cite: 3].
\end{itemize}


\section{Visualización del Camino de Datos para Instrucción Tipo R (add)}

Mediante la herramienta integrada \textbf{Tools $\rightarrow$ MIPS X-Ray}, para una instrucción Tipo R (como \texttt{add \$s0, \$a1, \$a2})[cite: 3]:
\begin{enumerate}
    \item \textbf{IF:} La instrucción es leída desde la memoria de instrucciones y el PC se incrementa ($PC + 4$)[cite: 3].
    \item \textbf{ID:} La unidad de control decodifica el \textit{opcode} y lee los valores contenidos en los registros fuente (\texttt{\$a1} y \texttt{\$a2}) desde el banco de registros[cite: 3].
    \item \textbf{EX:} La ALU procesa la suma utilizando ambos operandos[cite: 3].
    \item \textbf{MEM:} Etapa inactiva. No se realiza ninguna interacción con la memoria de datos[cite: 3].
    \item \textbf{WB:} El resultado generado por la ALU se escribe de forma definitiva en el registro destino (\texttt{\$s0})[cite: 3].
\end{enumerate}


\section{Visualización del Camino de Datos para Instrucción Tipo I (lw)}

Mediante \textbf{Tools $\rightarrow$ MIPS X-Ray}, para una instrucción Tipo I de carga de memoria (como \texttt{lw \$t4, 0(\$t3)})[cite: 3]:
\begin1
    \item \textbf{IF:} La instrucción se obtiene de la memoria apuntada por el PC y este se incrementa ($PC + 4$)[cite: 3].
    \item \textbf{ID:} Se decodifica el \textit{opcode}, se identifica el registro base (\texttt{\$t3}), el destino (\texttt{\$t4}) y el valor inmediato (offset), el cual se extiende en signo a 32 bits[cite: 3].
    \item \textbf{EX:} La ALU suma la dirección base con el inmediato extendido para calcular la dirección efectiva[cite: 3].
    \item \textbf{MEM:} Se activa completamente. El procesador accede a la memoria de datos usando la dirección calculada y extrae la palabra de 32 bits[cite: 3].
    \item \textbf{WB:} El dato leído de la memoria se escribe permanentemente en el registro destino (\texttt{\$t4})[cite: 3].
\end{enumerate}


\section{Justificación de los Algoritmos Seleccionados}

\begin{itemize}
    \item \textbf{Relevancia Conceptual y Recursividad:} Quicksort y Mergesort requieren una gestión recursiva explícita, sirviendo para poner en práctica la construcción manual de marcos de pila (\textit{stack frames}) y la preservación de registros según el MIPS ABI[cite: 3].
    \item \textbf{Análisis In-Place vs. Memoria Auxiliar:} Quicksort permite analizar un enfoque \textit{in-place} mediante intercambios (\textit{swaps}), mientras que Mergesort evalúa el costo de un búfer auxiliar en RAM (\texttt{temp}) y el tráfico masivo de instrucciones \texttt{lw}/\texttt{sw}[cite: 3].
    \item \textbf{Eficiencia Algorítmica:} Demuestra cómo la reducción de complejidad asintótica a $\mathcal{O}(n \log n)$ compensa el \textit{overhead} derivado del manejo de la pila (\texttt{addi \$sp, \$sp, -20})[cite: 3].
\end{itemize}


\section{Análisis y Discusión de los Resultados}

El estudio demuestra que la eficiencia en MIPS32 depende del equilibrio entre la complejidad teórica y la interacción con el camino de datos y la memoria[cite: 3]. Mergesort ofrece un tiempo $\mathcal{O}(n \log n)$ garantizado, pero su fase de mezcla en el búfer auxiliar sobrecarga el bus de datos con instrucciones de carga y almacenamiento[cite: 3]. Quicksort optimiza los accesos a RAM al operar \textit{in-place}, aunque está expuesto a peores casos $\mathcal{O}(n^2)$ según la elección del pivote[cite: 3].

Asimismo, la disciplina en el manejo de la pila (\texttt{addi \$sp, \$sp, -20}) resulta vital para prevenir corrupciones durante la recursión[cite: 3]. Las herramientas de depuración de MARS (paso a paso, panel de registros y MIPS X-Ray) permitieron validar el camino de datos (IF-ID-EX-MEM-WB) y confirmaron que la máxima optimización se logra combinando baja complejidad algorítmica con un uso intensivo del banco de registros internos[cite: 3].

\end{document}
